\chapter{Ricci Flow as a Gradient Flow}

\section{6.1 \quad Gradient of total scalar curvature and related functionals}

For a compact Riemannian manifold $(\mathcal M,g)$, define
\[
E(g)=\int_{\mathcal M}R\,dV.
\]

\begin{proposition}
\label{topping-gradient-flow-total-scalar-curvature-gradient}
\source{topping:page-0068,topping:page-0069}
\dcref{ch6:1.1}
\group{topping-gradient-flow}
The $L^2$-gradient of the total scalar curvature is
\[
\nabla E(g)=\frac{R}{2}g-\Ric.
\]
If the metric formally follows the upward gradient flow
\[
\frac{\partial g}{\partial t}=\frac{R}{2}g-\Ric,
\]
then
\[
\frac{\partial R}{\partial t}
=-\left(\frac{n}{2}-1\right)\Delta R+|\Ric|^2-\frac12R^2.
\tag{6.1.1}
\]
\end{proposition}

\begin{proof}
\sketch For $h=\partial_tg$, the first variation is
\[
\frac{d}{dt}E(g)=\int_{\mathcal M}
\left\langle\frac{R}{2}g-\Ric,h\right\rangle dV,
\]
after the divergence terms integrate to zero. Substitution of
$h=\frac{R}{2}g-\Ric$ into the scalar-curvature variation formula, together
with $\delta\Ric+\frac12dR=0$, gives the displayed evolution equation.
\end{proof}

For $n\geq3$, the leading term in this scalar-curvature equation has the
backward-heat sign, so the upward gradient equation is not generally a
forward well-posed evolution from arbitrary initial metrics.

\section{6.2 \quad The $\mathcal{F}$-functional}

For a fixed closed manifold $\mathcal M$, the $\mathcal F$-functional on a
metric $g$ and a smooth function $f:\mathcal M\to\mathbb R$ is
\[
\mathcal F(g,f)=\int_{\mathcal M}(R+|\nabla f|^2)e^{-f}\,dV.
\]

\begin{proposition}
\label{topping-gradient-flow-f-functional-first-variation}
\source{topping:page-0069,topping:page-0070}
\dcref{ch6:2.1}
\group{topping-gradient-flow}
For a variation satisfying $\partial_tg=h$ and $\partial_tf=k$,
\[
\frac{d}{dt}\mathcal F(g,f)
=\int_{\mathcal M}\langle-\Ric-\Hess(f),h\rangle e^{-f}\,dV
+\int_{\mathcal M}\left(2\Delta f-|\nabla f|^2+R\right)
\left(\frac12\tr h-k\right)e^{-f}\,dV.
\]
\end{proposition}

\begin{proof}
\sketch Differentiate $|\nabla f|^2$ to obtain
\[
\partial_t|\nabla f|^2=-h(\nabla f,\nabla f)+2\langle\nabla k,\nabla f\rangle.
\]
Combine this with the scalar-curvature and volume variation formulas, then
integrate the $\nabla k$, divergence, and Laplacian terms by parts against
$e^{-f}dV$.
\end{proof}

\section{6.3 \quad The heat operator and its conjugate}

For functions on $\mathcal M\times[0,T]$, set
\[
\Box=\frac{\partial}{\partial t}-\Delta,
\qquad
\Box^*=-\frac{\partial}{\partial t}-\Delta+R.
\]

\begin{proposition}
\label{topping-gradient-flow-conjugate-heat-operator}
\source{topping:page-0071}
\dcref{ch6:3.1}
\group{topping-gradient-flow}
If $g(t)$ is a Ricci flow on $[0,T]$ and $v,w$ are smooth functions on
$\mathcal M\times[0,T]$, then
\[
\int_0^T\!\int_{\mathcal M}(\Box v)w\,dV\,dt
=\left[\int_{\mathcal M}vw\,dV\right]_0^T
+\int_0^T\!\int_{\mathcal M}v(\Box^*w)\,dV\,dt.
\]
\end{proposition}

\begin{proof}
\sketch Differentiate $\int_{\mathcal M}vw\,dV$ using the Ricci-flow
volume variation, integrate the spatial Laplacian terms by parts, and collect
the resulting terms as $\Box v$ and $\Box^*w$.
\end{proof}

\begin{remark}
\label{topping-gradient-flow-conjugate-heat-special-case}
\source{topping:page-0071}
\dcref{ch6:3.2}
\group{topping-gradient-flow}
Taking $v\equiv1$ in the conjugacy identity gives
\[
\frac{d}{dt}\int_{\mathcal M}w\,dV=-\int_{\mathcal M}\Box^*w\,dV.
\]
\end{remark}

\section{6.4 \quad A gradient flow formulation}

Consider variations of $(g,f)$ preserving the weighted volume form
$e^{-f}dV$. The constraint and the resulting first variation are
\[
0=\partial_t(e^{-f}dV)=\left(\frac12\tr(\partial_tg)-\partial_tf\right)e^{-f}dV,
\qquad
\partial_tf=\frac12\tr(\partial_tg),
\]
and hence
\[
\frac{d}{dt}\mathcal F(g,f)=
\int_{\mathcal M}\left\langle-\Ric-\Hess(f),\partial_tg\right\rangle e^{-f}\,dV.
\]

The weighted-volume-preserving gradient system is
\begin{align}
\frac{\partial g}{\partial t}&=-2(\Ric+\Hess(f)), \tag{6.4.1}\\
\frac{\partial f}{\partial t}&=-R-\Delta f, \tag{6.4.2}\\
\frac{d}{dt}\mathcal F(g,f)&=2\int_{\mathcal M}|\Ric+\Hess(f)|^2e^{-f}\,dV\geq0. \tag{6.4.3}
\end{align}
With
\begin{equation}
\omega=e^{-f}dV, \tag{6.4.4}
\end{equation}
the form $\omega$ is constant in time, and $f$ is determined by $\omega$ and
$g$. The coupled system is related by time-dependent diffeomorphisms to
\begin{align}
\frac{\partial g}{\partial t}&=-2\Ric, \tag{6.4.5}\\
\frac{\partial f}{\partial t}&=-\Delta f+|\nabla f|^2-R. \tag{6.4.6}
\end{align}

Given $g_0$, short-time Ricci-flow existence supplies a smooth solution of
(6.4.5) on $[0,T]$ for some $T>0$. For this metric evolution, prescribe
smooth terminal data $f(T)$ and solve (6.4.6) backward. With
\begin{equation}
u(t)=e^{-f(t)}, \tag{6.4.7}
\end{equation}
the transformed function satisfies
\begin{equation}
\Box^*u=0, \tag{6.4.8}
\end{equation}
and the conjugate heat equation has a positive solution backward to $t=0$.

Let $X=-\nabla f$ and let $\psi_t$ be its flow with $\psi_0=\id$.
Set $\hat g=\psi_t^*g$ and $\hat f=f\circ\psi_t$. Since
$\Lie_Xg=-2\Hess_g(f)$, the pullback equations are
\begin{align}
\frac{\partial\hat g}{\partial t}
&=-2\bigl(\Ric(\hat g)+\Hess_{\hat g}(\hat f)\bigr), \tag{6.4.9}\\
\frac{\partial\hat f}{\partial t}
&=-\Delta_{\hat g}\hat f-R_{\hat g}. \tag{6.4.10}
\end{align}
Thus solutions of (6.4.5)--(6.4.6) produce solutions of (6.4.1)--(6.4.2),
and conversely by the inverse diffeomorphism. Pullback by a common
diffeomorphism leaves $\mathcal F$ invariant, so $\mathcal F$ is
nondecreasing along (6.4.5)--(6.4.6).

\begin{theorem}
\label{topping-gradient-flow-gradient-flow-formulation}
\source{topping:page-0074}
\uses{topping-gradient-flow-f-functional-first-variation, topping-gradient-flow-conjugate-heat-operator, topping:variation-volume-form}
\dcref{ch6:4.1}
\group{topping-gradient-flow}

For every initial metric $g_0$ on a closed manifold $\mathcal M$, there are
$T>0$ and a positive $n$-form $\omega$, generally neither arbitrary nor
unique, such that on $[0,T]$ the Ricci flow starting at $g_0$ is, up to a
time-dependent diffeomorphism, an upward gradient flow of
\[
\mathcal F^\omega(g)=\int_{\mathcal M}(R+|\nabla f|^2)\,\omega,
\qquad \omega=e^{-f}dV,
\]
with respect to
\[
\langle g,h\rangle_\omega=\frac12\int_{\mathcal M}\langle g,h\rangle\,\omega.
\]
\end{theorem}

\begin{proof}
\sketch Solve the Ricci flow and the conjugate heat equation with terminal
data, then pull back by the flow of $-\nabla f$. The pulled-back pair obeys
the weighted-volume-preserving gradient system, and the first variation
identifies it as the upward gradient flow for $\mathcal F^\omega$.
\end{proof}

The interval endpoint $T$ can be chosen within the existence interval of the
Ricci flow, but the corresponding $\omega$ depends on that choice. The
gradient flow therefore represents Ricci flow modulo a time-dependent family
of diffeomorphisms.

\begin{remark}
\label{topping-gradient-flow-steady-soliton-criterion}
\source{topping:page-0074}
\dcref{ch6:4.3}
\group{topping-gradient-flow}
If $d\mathcal F/dt=0$ along (6.4.1)--(6.4.2), then
$\Ric+\Hess(f)=0$. On a compact manifold the associated Ricci flow is Ricci
flat, so it is a steady gradient soliton.
\end{remark}

\section{6.5 \quad The classical entropy}

For a positive solution $u$ of the conjugate heat equation, define the
Boltzmann-Shannon entropy
\[
N=-\int_{\mathcal M}u\ln u\,dV.
\]
Using $\Box^*u=0$ and the Ricci-flow volume variation,
\begin{align}
-\frac{dN}{dt}
&=\int_{\mathcal M}\left[(1+\ln u)\frac{\partial u}{\partial t}-Ru\ln u\right]dV \\
&=-\int_{\mathcal M}\left[(1+\ln u)(\Delta u-Ru)+Ru\ln u\right]dV \\
&=\int_{\mathcal M}\left(\frac{|\nabla u|^2}{u}+Ru\right)dV
=\mathcal F. \tag{6.5.1}
\end{align}
Define the renormalized entropy by
\begin{equation}
\widetilde N=N-\frac n2\left(1+\ln\bigl(4\pi(T-t)\bigr)\right). \tag{6.5.2}
\end{equation}
For the conjugate heat kernel with terminal delta data, this normalization
satisfies $\widetilde N\to0$ as $t\uparrow T$.

\begin{lemma}
\label{topping-gradient-flow-classical-entropy-monotone}
\source{topping:page-0076}
\uses{topping-gradient-flow-gradient-flow-formulation}
\dcref{ch6:5.1}
\group{topping-gradient-flow}

If $g(t)$ is a Ricci flow on a closed manifold $M$ for $t\in[0,T]$ and
$u:M\times[0,T]\to(0,\infty)$ satisfies $\Box^*u=0$, then the functional
$\widetilde N$ in (6.5.2) is weakly increasing in $t$.
\end{lemma}

\begin{proof}
\sketch Since
\[
-\frac{d\widetilde N}{dt}=\mathcal F-\frac{n}{2(T-t)},
\]
it suffices to show
\begin{equation}
\mathcal F\leq\frac{n}{2(T-t)}. \tag{6.5.3}
\end{equation}
The gradient-flow formulation and the trace estimate give
\begin{align}
\frac{d\mathcal F}{dt}
&=2\int_{\mathcal M}|\Ric+\Hess(f)|^2e^{-f}\,dV \\
&\geq\frac2n\int_{\mathcal M}|R+\Delta f|^2e^{-f}\,dV \\
&\geq\frac2n\left(\int_{\mathcal M}(R+\Delta f)e^{-f}\,dV\right)^2
=\frac2n\mathcal F^2, \tag{6.5.4}
\end{align}
where $f=-\ln u$. The first inequality uses the tensor trace bound, the
second uses Cauchy--Schwarz, and integration by parts identifies the final
integral with $\mathcal F$. If (6.5.3) failed, (6.5.4) would force $\mathcal F$
to blow up before $T$, contradicting its existence on $[0,T]$.
\end{proof}

\section{6.6 \quad The zeroth eigenvalue of \texorpdfstring{$-4\Delta + R$}{-4 Delta + R}}

Writing $\phi=e^{-f/2}$ gives
\[
\mathcal F(g,\phi)=\int_{\mathcal M}(4|\nabla\phi|^2+R\phi^2)\,dV.
\]
The zeroth eigenvalue of $-4\Delta+R$ is therefore
\[
\lambda(g)=\inf\left\{\mathcal F(g,\phi)\,\middle|\,
\int_{\mathcal M}\phi^2\,dV=1\right\}.
\]
For smooth $g$ on a closed manifold, the direct method and elliptic regularity
(see [15, \S8.12]) give a smooth minimizer $\phi_{\min}$; after changing its
sign, it is positive by the Harnack inequality (see [15, Theorem 8.20]).
Positivity implies uniqueness up to sign, since a nontrivial linear
combination of two distinct minimizers would otherwise have a zero. Writing
$\phi_{\min}=e^{-f_{\min}/2}$ yields
\begin{equation}
\lambda=\min\left\{\mathcal F(g,f)\,\middle|\,
\int_{\mathcal M}e^{-f}\,dV=1\right\}. \tag{6.6.1}
\end{equation}

\begin{theorem}
\label{topping-gradient-flow-z0-eigenvalue-monotone}
\source{topping:page-0077,topping:page-0078}
\uses{topping-gradient-flow-gradient-flow-formulation}
\dcref{ch6:6.1}
\group{topping-gradient-flow}

If $g(t)$ is a Ricci flow on a closed manifold $\mathcal M$ for
$t\in[0,T]$, then $\lambda(g(t))$ is weakly increasing on $[0,T]$.
\end{theorem}

\begin{proof}
\sketch Let $0\leq a<b\leq T$, and let $f_b$ minimize (6.6.1) for $g(b)$.
Solve (6.4.6) backward on $[a,b]$ with terminal value $f(b)=f_b$. The
monotonicity of $\mathcal F$ gives
\[
\lambda(g(a))\leq\mathcal F(g(a),f(a))
\leq\mathcal F(g(b),f(b))=\lambda(g(b)).
\]
\end{proof}
